\subsection{Money}
\label{sec:11.4.3}
Against a state operator, Congress's ordinary lever is a condition on its appropriation. The abduction and forcible transfer pipeline has been funded so that the condition seldom bites.\note{
Since 2025 most of the money has come by reconciliation and stays available through fiscal 2029: \$45 billion for detention capacity and nearly \$30 billion more for enforcement \autocite{obbba2025}, and roughly \$70 billion more for ICE and CBP \autocite{secureamerica2026,gao2026detention}. For fiscal 2026 ICE received no full-year appropriation, and when DHS appropriations lapsed on 14 February 2026 it went on operating on the reconciliation money, which stayed available throughout \autocite{gao2026detention}.

Annual appropriation, the instrument Boland used, now reaches little of the money. Multi-year lump sums without line items also give a funds limitation less to attach to. Amnesty's letter asks Congress to rescind ICE's “current funding” \autocite{AmnestyUSA2026Letter}. Most of that funding is now multi-year balances, which the Abolish ICE Act's rescission reaches and next year's appropriations bill does not.

Appropriations can also run the other way, and the record shows how. From fiscal 2010 through fiscal 2016, DHS appropriations acts provided that ICE's funding “shall maintain a level of not less than” 33,400, later 34,000, detention beds \autocite{dhsapprop2010,Geng2018BedQuota}. The text required beds, not detainees, but ICE tracked its compliance by average daily population \autocite{Geng2018BedQuota}. A capacity floor therefore worked as an occupancy target: a rate the funder fixed in advance and the operator met by finding people.

The same 2010 act carried the other kind of condition too. None of the funds under its ICE salaries heading could continue a 287(g) delegation if the DHS Inspector General found that the delegation's terms had been violated \autocite{dhsapprop2010}. The finding cut off the money by itself, with no decision by ICE in between, which is the shape of a review that binds.
}
